% Generic A0 portrait draft. Adapt dimensions/orientation to the approved goal.
\documentclass{ltjsarticle}
\usepackage[paperwidth=841mm,paperheight=1189mm,margin=38mm]{geometry}
\usepackage{graphicx,booktabs,xcolor,tikz}
\usepackage[hidelinks]{hyperref}
\definecolor{accent}{HTML}{174C65}
\pagestyle{empty}
\setlength{\parindent}{0pt}
\setlength{\parskip}{12mm}
\newcommand{\paneltitle}[1]{{\color{accent}\fontsize{48}{60}\selectfont\bfseries #1}\par\vspace{5mm}}
\begin{document}
\fontsize{30}{44}\selectfont
{\color{accent}\fontsize{76}{92}\selectfont\bfseries ポスタータイトル\par}
{\fontsize{40}{54}\selectfont 最も伝えたいメッセージを一文で示す。\par}
% Insert real author/affiliation information only when supplied.
\vspace{15mm}
{\color{accent}\rule{\linewidth}{2mm}}
\vspace{18mm}
\begin{minipage}[t][0.72\textheight][t]{0.46\linewidth}
\raggedright
\paneltitle{背景と問い}
対象とする課題と、なぜ重要なのかを簡潔に説明する。\par
\vspace{18mm}
{\color{accent}\rule{\linewidth}{0.5mm}}\par
読者が理解すべき背景を絞り、答えたい問いを明確にする。\par
\vfill
\paneltitle{方法}
対象、手順、評価方法を図と短い説明で示す。\par
\vspace{12mm}
\begin{tikzpicture}[every node/.style={draw=accent,rounded corners,minimum width=90mm,minimum height=50mm,font=\fontsize{28}{36}\selectfont}]
\node (a) at (0,0) {入力};
\node (b) at (14,0) {分析};
\draw[->,line width=2mm,accent] (a) -- (b);
\end{tikzpicture}
\par\vspace{18mm}
図の要素は実際の手順に置き換え、再現に必要な条件を短く補足する。\par
\vfill
\paneltitle{根拠と条件}
データの出典、評価条件、比較対象を示す。未実施の検証や未確認の数値を補わない。\par
\end{minipage}\hfill
\begin{minipage}[t][0.72\textheight][t]{0.46\linewidth}
\raggedright
\paneltitle{主要な結果}
確認済みの結果を大きな図表で示す。\par
\vspace{18mm}
\begingroup\setlength{\fboxsep}{12mm}\color{accent}
\fbox{\begin{minipage}[c][240mm][c]{0.9\linewidth}
\centering\fontsize{40}{54}\selectfont 主図の配置領域\par
\fontsize{28}{40}\selectfont 実際の成果に基づく図へ置き換える。\par
\end{minipage}}\endgroup\par
\vspace{12mm}
図から読み取れることと、支持される結論を対応させる。\par
\vfill
\paneltitle{結論と限界}
問いへの答え、根拠、適用範囲を短くまとめる。\par
\vspace{18mm}
不確実性や未解決の問いを明示し、中心的なメッセージを持ち帰れる構成にする。\par
\end{minipage}
\vfill
{\color{accent}\rule{\linewidth}{1mm}}
{\fontsize{24}{34}\selectfont 参考文献・連絡先：確認済みの情報のみを記載する。\par}
\end{document}
